\section{Bidirectional derivation of the electronic exponent from the enriched route}
\label{sec:cierre-electron-two-way-rev11}
\index{electron!bidirectional reader}
\index{calendar of order 108! electronic exponent}

The enriched route of the central point publishes two complementary representations
of its memory: the complete tritic word
\(\gamma_{108}\in\mathbb F_3^{108}\) and the signed dodecaphase record
\(\tau_{12}\in\{-1,0,+1\}^{12}\). The two readers defined below
are computed exclusively from those representations, the calendar, and the
HMT constants already constructed.

\subsection{Reader of displacements on the cyclic calendar of order 108}

Let \(S\) be the cyclic shift on \(\mathbb Z/108\mathbb Z\). For
\(k\in\mathbb Z/108\mathbb Z\), define
\begin{equation}
 D_k(\Gamma):=
 \#\{t\in\mathbb Z/108\mathbb Z:\gamma_t\ne\gamma_{t+k}\}
 =\lVert(I-S^k)\gamma\rVert_0.
 \label{eq:lector-dk-ct108-rev11}
\end{equation}
The integer triple of the reader is
\begin{equation}
 K_D(\Gamma)=
 \left(D_{27}+D_{36},\ D_1,\ \frac{D_4-D_3}{2}\right)
 =:(A_D,B_D,C_D),
 \label{eq:KD-rev11}
\end{equation}
and its scalar evaluation is
\begin{equation}
 \kappa_D(\Gamma)=A_D-\alpha_{\rm HMT}B_D+\Delta_4 C_D.
 \label{eq:kappaD-rev11}
\end{equation}
The shifts (27) and (36) are the chronological lifts of
\(90^\circ\) and \(120^\circ\) because nine discrete steps represent
\(30^\circ\). The term \(D_1\) measures the immediate local boundary, and the
difference \(D_4-D_3\) compares the two genealogies before their height
quotient. In the domain of the 324 routes,
\(\gamma_{t+54}=\gamma_t\) holds; all the \(D_k\) are even and \(C_D\) is an integer.

\subsection{Dodecaphasic correlation difference reader}

For the word \(\tau=(\tau_r)_{r\in\mathbb Z/12\mathbb Z}\), let
\[
 C_k^{\rm corr}(\tau)=\sum_{r=0}^{11}\tau_r\tau_{r+k},
 \qquad
 A_\ell(\tau)=\sum_{r=0}^{11}
 \left(\sum_{j=0}^{\ell-1}\tau_{r+j}\right)^2.
\]
The primitive genealogical obstruction and the antipodal memory are
\begin{align}
 \Omega_{90\mid120}(\tau)
 &=4A_3-3A_4
 =-2C_1^{\rm corr}-4C_2^{\rm corr}-6C_3^{\rm corr},
 \label{eq:omega-two-way-rev11}\\
 P_6(\tau)&=\sum_{r=0}^{5}\tau_r\tau_{r+6}.
 \label{eq:p6-two-way-rev11}
\end{align}
The symbol \(P_6\) prevents the colision with the short cycle of the corredor of
constantes. The second triple and its evaluation are
\begin{equation}
 K_\Omega(\Gamma)=\bigl(\Omega_{90\mid120}(\tau_\Gamma),54,
 P_6(\tau_\Gamma)\bigr),
 \qquad
 \kappa_\Omega=\Omega_{90\mid120}-54\alpha_{\rm HMT}+P_6\Delta_4.
 \label{eq:kappa-omega-rev11}
\end{equation}

\HMTClaim{MD-R-ELECTRON-OMEGA-PRIMITIVE}
\begin{lemma}[Linear integral primitive uniqueness]
\label{lem:unicidad-omega-rev11}
Let \(L=uA_3+vA_4\), with \(u,v\in\mathbb Z\), be a contrast that vanishes
when the genealogies descend to the same height,
\(A_3=3a\), \(A_4=4a\). Then
\((u,v)=n(4,-3)\). Once the sign \(90-120\) has been fixed,
\(\Omega_{90\mid120}=4A_3-3A_4\) is the primitive generator.
\end{lemma}

\begin{proof}
Vanishing for every \(a\) requires \(3u+4v=0\). Since
\(\gcd(3,4)=1\), the integer solutions are
\((u,v)=n(4,-3)\). Primitivity corresponds to \(|n|=1\).
\end{proof}

The lemma selects \(\Omega\) within the class of integral linear contrasts
of the energies \(A_3,A_4\). Its uniqueness domain is that functional
class; nonlinear invariants belong to a distinct selection problem.

\subsection{Evaluation of the central route}

Before metrological comparison, the electronic route \(\Gamma_e=(5,5,++)\) provides,
\[
 (D_1,D_3,D_4,D_{27},D_{36})=(54,56,68,48,32),
\]
and its signed dodecaphonic register is
\[
 \tau_e=(+,+,+,-,-,-,+,+,+,-,-,-).
\]
Therefore,
\[
 K_D(\Gamma_e)=(48+32,54,(68-56)/2)=(80,54,6),
\]
\[
 K_\Omega(\Gamma_e)=(80,54,6).
\]

\HMTClaim{MD-R-ELECTRON-TWOWAY-REDUCTION}
\begin{theorem}[Common Electronic Reduction of Bidirectional Readers]
\label{thm:reduccion-electron-two-way-rev11}
The readers \(K_D\) and \(K_\Omega\) coincide on the central route and
determine
\begin{equation}
 \kappa_e=80-54\alpha_{\rm HMT}+6\Delta_4.
 \label{eq:kappa-electron-demostrado-rev11}
\end{equation}
\end{theorem}

\begin{proof}
The two preceding evaluations yield the same integer triple. Substitution
of its components into
\(A-\alpha_{\rm HMT}B+\Delta_4C\) proves
\eqref{eq:kappa-electron-demostrado-rev11}. Every quantity used in this
reduction belongs to the route, to its two records, or to the previously
constructed HMT constants.
\end{proof}

\paragraph{Subsequent realization on the energy chart.}
With \(R_{\rm act}=H_5/\etaRet\), the positive action on the baseline
\(\mathfrak e_0\) yields
\begin{equation}
 \mathfrak e_e^\beta
 =\mathfrak e_0R_{\rm act}^{\kappa_e}.
 \label{eq:masa-electron-adimensional-two-way-rev11}
\end{equation}
After applying the current chart to MeV is obtained
\begin{equation}
 \varsigma_{u_E}(\mathfrak e_e^\beta)\big|_{u_E=1\,\mathrm{MeV}}
 =0.5109989506900008086082571648\ldots\ \mathrm{MeV}.
 \label{eq:masa-electron-two-way-rev11}
\end{equation}
The first equality belongs to the internal theorem conditioned by the basal and
the action line; the second is its dimensional realization. The external
value is reserved for metrological comparison.

The integer (80) follows from \(D_{27}+D_{36}\) and, in the dodecaphase reader,
from \(\Omega_{90\mid120}\). Its coincidence with other internal cardinals is
a compatibility control and does not replace this operator provenance. The
orientation route \((--)\) is a temporal translation of route \((++)\)
and preserves the same profile, in accordance with mass parity under conjugation.

\subsection{Separation of readers and return scales}

Numerical coincidence of internal components is subordinate to their
type. At the central point, the objects involved include
\[
 R_{\rm atlas}(5,5)=(24,0,12,0,-12),
 \qquad
 L_{\kappa_0}(\Gamma_e^{++})=(24,0,0,0,-12),
\]
\[
 \begin{gathered}
 k_{\Delta_4}^{\rm dod},k_{\Delta_4}^{\rm pro}:
 \mathscr R_{324}\longrightarrow\mathbb Z,\\
 k_{\Delta_4}^{\rm dod}(\Gamma_e^{++})=12,
 \qquad
 k_{\Delta_4}^{\rm pro}(\Gamma_e^{++})=4,\\
 K_D(\Gamma_e^{++})=K_\Omega(\Gamma_e^{++})=(80,54,6).
 \end{gathered}
\]
Here \(\mathscr R_{324}\) denotes the domain of the 324 oriented routes.
The first is a static family signature; the second, a genealogy
extractor; the third and fourth terms are torsional readers defined,
respectively, by dodecaphasic reading and prospective prolongation;
the final object is the pair of bidirectional readers. No partial equality
identifies those morphisms or permits transferring a coefficient between
readers.

Four temporal scales are also distinguished: \(27\) is the first return of class, orbit, and cell; \(54=D_1(\Gamma_e)\) is the observable kinematic return within the cycle; \(108\) closes the calendar; and \(216\) closes the oriented kinematic state in the extension. This typing keeps positional return, remark, and oriented holonomy separate.

\paragraph{Causal chain of the electronic outcome}
The route and its two registers determine \(K_D\), \(K_\Omega\), their common
reduction, and the evaluation of \(\kappa_e\). The realization in MeV subsequently applies
the energy chart; external mass enters only in the comparison. The pair
\(\mathbf K=(K_D,K_\Omega)\) is the complete operator output of HMT on
each route. In a noncentral fiber, its representation by a single physical
scalar requires an additional morphism between the fiber operator and the
dimensional chart; this requirement neither modifies the internal law nor authorizes
the selection of a cell by metrological proximity.
